% page 467 du fac-simile (image p-466 du scan)
% bloc 162.6 x 237.7 mm ; marge gauche 39.4 mm
\pgc{17.780mm}{0.000mm}{
\l{\cel{51}459}
\l{}
\l{\sou{presbiterismo}. (\sou{religio}).Doktrino di ula sekto Kristana qua ne}
\l{\cel{3}agnoskas la episkopi, e grantas la direkto di la eklezio a}
\l{\cel{3}la sacerdoti o pastori. - DEFIRS.}
\l{}
\l{\sou{presbiterio}. (\sou{religio)}. Asemblitaro ek la personi qui adminis-}
\l{\cel{3}tras la proprietaji e la revenui di la parokio. - DEFIRS.}
\l{}
\l{\sou{preske}. \textquotedbl{}Poko mankas pri ke lo esez tale\textquotedbl{}.}
\l{}
\l{\sou{preskriptar}. (\sou{trans., ulo ad ulu)}.Imperar per formulo preciza.}
\l{\cel{3}EFIS.}
\l{}
\l{\sou{preskritar}. (\sou{trans.}) (\sou{yurocienco}).Liberigar de debajo, de perse-}
\l{\cel{3}qui judiciala, de servitudo, per la fakto ke quanto determi-}
\l{\cel{3}nita de tempo fluis de lore. - FIS.}
\l{}
\l{\sou{prestar}. (\sou{trans., ad)}.\cel{2}Igar (ulo) disponebla, posedata, da ulu}
\l{\cel{3}dum ula quanto de tempo, pos qua lu devos retrodonar olu.-}
\l{\cel{3}FIS.}
\l{}
\l{\sou{prestaciono.}\cel{2}Laboro kelka-dia (redemtebla per pekunio) postul-}
\l{\cel{3}ebla de la habitanti di komono, por la konservo en bona stan-}
\l{\cel{3}do di la voyi publika qui apartenas a la komono. - DEFIS.}
\l{}
\l{\sou{prestijo}. To quo frapas, astonas, surprizas per sua brilo, e c.}
\l{\cel{3}Autoritato od influo qua naskas de la aspekto extera, ed}
\l{\cel{3}efikas a la opiniono ed a la imagino-fakultato di ula per-}
\l{\cel{3}soni.-DEFIRS.}
\l{}
\l{\sou{pretend}ar. (\sou{trans.}) Persequar publike ulo pri quo onu dicas}
\l{\cel{3}esar yuroza. - DEFIRS.}
\l{}
\l{\sou{preter}. Pasante apud (ulu od ulo) ed irante plu ad-avane kam}
\l{\cel{3}(li). L.}
\l{}
\l{\sou{preterito}. (\sou{gram.}) Tempo di verbo, qua expresas la pasinto. -}
\l{\cel{3}DEFIS.}
\l{}
\l{\sou{pretextar}.\cel{2}(\sou{trans.}) Alegar motivo qua aspektas exakta, por ce-}
\l{\cel{3}lar la motivo reala di ago. - EFIS.}
\l{}
\l{\sou{pretoro}. (en\cel{1}\sou{Roma, antique)}. Oficiero judiciala qua direktis}
\l{\cel{3}la operaci militala, administris provinco, e c. - DEFIS.}
\l{}
\l{\sou{prevarikar}. (\sou{netrans.}) Deviacar de la konduto yuro-konforma.}
\l{\cel{3}- FIS.}
\l{}
\l{\sou{preventar.}\cel{2}(\sou{trans}.)Impedar la evento, la existesko. - EFIS.}
\l{}
\l{\sou{prevosto}. I. Oficiero di qua la tasko konsistas ye judiciar}
\l{\cel{3}pri la kazi kriminala en la armei. - II. Chefo di la kanoni-}
\l{\cel{3}karo en kirko kolegiala, superioro en ula ordeni religiala.}
\l{\cel{3}- EFI.}
}
